\documentclass[main.tex]{subfiles}
\input{tex/b46/frankl/edition-setup.tex}
\input{tex/b46/frankl/statements.tex}
\input{tex/b46/frankl/proof-display.tex}
\input{tex/b46/frankl/triple-statement.tex}
\providecommand{\ProofRefStack}[1]{%
  \ensuremath{\begin{gathered}[t]#1\end{gathered}}}
\title{Warum ein Element in mindestens\\der Hälfte vorkommt\\[.5em]\large Frankls Vermutung bei kleinen Mengengliedern\\Beweistabellen zu Band 46}
\author{Martin Kunze}
\date{}
\hypersetup{pdftitle={Band 46 - Frankls Spezialfälle - Beweistabellen},pdfauthor={Martin Kunze}}
\begin{document}
\hypersetup{pageanchor=false}
\maketitle
\hypersetup{pageanchor=true,hypertexnames=false}
\hypertarget{frankl.proofs}{}
\ProofWideReasonsNearFormula
\input{tex/b46/frankl/proofs.tex}
\clearpage
\renewcommand{\FormulaBandID}{46E}
\setcounter{chapter}{0}
\setcounter{section}{0}
% Überschrift, Voraussetzungen und Formel der neuen Aussagen zusammenhalten.
\NewCommandCopy\FranklTripleOriginalThm\FormulaThmDeltaK
\NewCommandCopy\FranklTripleOriginalDef\FormulaDefDeltaK
\RenewDocumentCommand{\FormulaThmDeltaK}{o m m +m}{%
  \par\Needspace{12\baselineskip}%
  \IfNoValueTF{#1}{\FranklTripleOriginalThm{#2}{#3}{#4}}{%
    \FranklTripleOriginalThm[#1]{#2}{#3}{#4}}%
}
\RenewDocumentCommand{\FormulaDefDeltaK}{o m m +m}{%
  \par\Needspace{12\baselineskip}%
  \IfNoValueTF{#1}{\FranklTripleOriginalDef{#2}{#3}{#4}}{%
    \FranklTripleOriginalDef[#1]{#2}{#3}{#4}}%
}
\input{tex/b46/frankl/triple-notation.tex}
\input{tex/b46/frankl/triple-proofs.tex}
\end{document}
